% ---------------------------------------------------------------------
%  PHỤ LỤC
% ---------------------------------------------------------------------
\chapter{Phụ lục}
\markboth{Phụ lục}{}
% Phu luc nam sau \backmatter nen chuong khong duoc danh so va bo dem bang, hinh, doan ma
% khong ve 0: de nguyen thi bang o day mang so tiep theo cua Chuong 5. Danh so rieng PL.1,
% PL.2... (ca ban \theH... de lien ket trong PDF khong trung voi bang cua Chuong 5).
\setcounter{table}{0}\renewcommand{\thetable}{PL.\arabic{table}}\renewcommand{\theHtable}{PL.\arabic{table}}
\setcounter{figure}{0}\renewcommand{\thefigure}{PL.\arabic{figure}}\renewcommand{\theHfigure}{PL.\arabic{figure}}
\setcounter{lstlisting}{0}\renewcommand{\thelstlisting}{PL.\arabic{lstlisting}}\renewcommand{\theHlstlisting}{PL.\arabic{lstlisting}}

\section*{A. Hướng dẫn cài đặt và chạy hệ thống}
\phantomsection
\addcontentsline{toc}{section}{A. Hướng dẫn cài đặt và chạy hệ thống}

Cách nhanh nhất là dùng Docker, vì hệ thống cần PostgreSQL có phần mở rộng
pgvector.

\begin{lstlisting}[language=bash, caption={Dựng toàn bộ hệ thống bằng Docker Compose.}]
# 1. Get the source code
git clone <repository-url> FAMS
cd FAMS

# 2. Create the environment file, then edit POSTGRES_PASSWORD and SECRET_KEY
copy .env.example .env

# 3. Build and start database, backend and web frontend
docker compose up --build

# 4. Load sample data (units, staff accounts, classes, schedules)
docker compose exec backend python seed.py

# 5. Open the web application
#    http://localhost:5173
\end{lstlisting}

Chạy từng phần khi phát triển:

\begin{lstlisting}[language=bash, caption={Chạy máy chủ và ứng dụng khi phát triển.}]
# Backend: needs Python 3.10 and a running PostgreSQL with pgvector
cd backend
python -m venv venv && venv\Scripts\activate
pip install -r requirements.txt
copy .env.example .env
alembic upgrade head
python seed.py
uvicorn app.main:app --reload --port 8001
# API documentation: http://localhost:8001/docs

# Flutter app on the web
cd mobile
flutter pub get
flutter run -d chrome --web-port 5174 --dart-define=API_BASE_URL=http://localhost:8001

# Automated tests
cd backend && pytest
cd mobile && flutter test
\end{lstlisting}

\begin{luuy}
Lần chạy đầu tiên, thư viện InsightFace tự tải bộ mô hình \code{buffalo\_l}
nặng khoảng 300\,MB, nên cần mạng và một khoảng chờ. Mô hình chống giả mạo
phải tải riêng và đặt vào \code{backend/models/}; nếu thiếu tệp này, hệ thống
vẫn chạy nhưng bỏ qua bước kiểm tra chống giả mạo và ghi một dòng cảnh báo.
Trình duyệt chỉ cho phép dùng camera trên HTTPS hoặc \code{localhost}.
\end{luuy}

\section*{B. Tài khoản mẫu để chấm và chạy thử}
\phantomsection
\addcontentsline{toc}{section}{B. Tài khoản mẫu để chạy thử}

Các tài khoản dưới đây do \code{seed.py} tạo ra, dùng để hội đồng chạy thử
từng vai trò. Mật khẩu chỉ dành cho môi trường thử nghiệm.

\begin{table}[H]
\centering
\caption{Tài khoản mẫu theo từng vai trò.}
\label{tab:taikhoanmau}
\small
\begin{tabularx}{\textwidth}{@{}>{\raggedright\arraybackslash}p{6.0cm} >{\raggedright\arraybackslash}p{3.0cm} X@{}}
\toprule
Tài khoản & Mật khẩu & Vai trò \\
\midrule
\code{admin@fams.edu.vn}      & \code{admin123}  & Quản trị viên \\
\code{pdt01@fams.edu.vn}      & \code{pdt12345}  & Phòng Đào tạo \\
\code{ctsv01@fams.edu.vn}     & \code{ctsv1234}  & Phòng Công tác sinh viên \\
\code{khoa.cntt@fams.edu.vn}  & \code{khoa1234}  & Văn phòng Khoa CNTT \\
\code{ldk.cntt@fams.edu.vn}   & \code{ldk12345}  & Lãnh đạo Khoa CNTT \\
\code{ban01@fams.edu.vn}      & \code{ban12345}  & Ban ngành \\
\code{gv01@fams.edu.vn}       & \code{gv123}     & Giảng viên \\
\code{sv01@fams.edu.vn}       & \code{sv123}     & Sinh viên \\
\code{ph01@fams.edu.vn}       & \code{ph123}     & Phụ huynh của sv01 \\
\bottomrule
\end{tabularx}
\end{table}

{\small
\begin{xltabular}{\textwidth}{@{}>{\raggedright\arraybackslash\leavevmode}p{2.6cm} >{\raggedright\arraybackslash\leavevmode}X@{}}
\caption{Giải thích chi tiết Bảng~\ref{tab:taikhoanmau}: nên thử gì với từng tài khoản.}
\label{tab:taikhoanmau-gt}\\
\toprule
Tài khoản & Nên thử \\
\midrule
\endfirsthead
\toprule
Tài khoản & Nên thử \\
\midrule
\endhead
\bottomrule
\endfoot
\code{admin} & Menu đầy đủ, nhật ký hệ thống toàn trường, đặt lại mật khẩu cho mọi tài khoản \\
\code{pdt01} & Tạo lớp học phần, xếp lịch vào phòng học và thử xếp trùng phòng để thấy bị chặn \\
\code{ctsv01} & Duyệt khuôn mặt, ghi nhận đồng ý theo bản giấy, duyệt nghỉ có phép, đặt lại mật khẩu sinh viên và phụ huynh \\
\code{khoa.cntt} & Tạo lớp và lịch trong Khoa CNTT; không có nút duyệt nghỉ có phép \\
\code{ldk.cntt} & Duyệt nghỉ có phép, xem nhật ký của Khoa; không có nút sửa lớp \\
\code{ban01} & Chỉ thấy thống kê tổng hợp, không thấy tên hay ảnh sinh viên \\
\code{gv01} & Mở phiên, quét camera cả lớp, chiếu mã QR, xác nhận tay và sửa điểm danh kèm lý do \\
\code{sv01} & Đồng ý sinh trắc học, đăng ký khuôn mặt, tự điểm danh bằng link trong mã QR \\
\code{ph01} & Xem chuyên cần của con, nhận thông báo mỗi lần con được điểm danh \\
Lưu ý & Các mật khẩu trên chỉ dành cho môi trường thử; bản chạy công khai đã đổi toàn bộ mật khẩu \\
\end{xltabular}}

\section*{C. Công cụ sinh số liệu cho báo cáo}
\phantomsection
\addcontentsline{toc}{section}{C. Công cụ sinh số liệu cho báo cáo}

Mọi con số về quy mô hệ thống trong báo cáo này không gõ tay mà được quét
trực tiếp từ kho mã nguồn bằng một công cụ nhỏ viết bằng Python.

\begin{lstlisting}[language=bash, caption={Sinh lại số liệu của báo cáo.}]
# Regenerate every number, table and code excerpt used by the report
python baocao/tools/quet-du-an.py

# Check only: exits with code 1 when the report data is out of date
python baocao/tools/quet-du-an.py --kiem-tra

# Static check of the LaTeX sources before compiling
python baocao/tools/kiem-tra-tex.py

# Measure recognition timings and write them into the report table
cd backend && python tools/do_hieu_nang.py --anh-mot-mat <image> --ghi
\end{lstlisting}

Công cụ đọc mã nguồn và sinh ra bốn tệp bảng cùng một tệp định nghĩa các lệnh
số liệu, đặt trong thư mục \code{baocao/auto/}. Từ đó, số bảng dữ liệu, số
endpoint, số ca kiểm thử, các ngưỡng cấu hình, các đoạn mã in trong
Chương~\ref{ch:caidat} và mã commit của bản được mô tả đều lấy thẳng từ hệ
thống thật. Công cụ cũng so sánh với lần quét trước và ghi ra tệp
\code{auto/thay-doi.md} liệt kê những gì đã đổi, kèm gợi ý mục nào trong báo
cáo cần soát lại.

\begin{ghinho}
Báo cáo này mô tả bản mã nguồn ở nhánh \code{\FnhanhGit}, commit
\code{\FcommitCuoi}, ngày \FngayCommit. Số liệu được quét lần gần nhất ngày
\FngayQuet.
\end{ghinho}

\section*{D. Bảng tự kiểm trước khi nộp}
\phantomsection
\addcontentsline{toc}{section}{D. Bảng tự kiểm trước khi nộp}

\begin{tukiem}
  \item Đã chạy \code{python baocao/tools/quet-du-an.py} lần cuối, số liệu
        trong báo cáo khớp với mã nguồn nộp kèm.
  \item Đã tìm chuỗi \emph{CẦN BỔ SUNG} trong tệp PDF và không còn chỗ nào
        sót lại.
  \item Đã điền họ tên, mã số sinh viên và lớp ở đầu tệp \code{do-an.tex}.
  \item Đã chèn ảnh chụp màn hình thật vào \code{figs/} và mở phần hình ở
        Chương~\ref{ch:caidat}.
  \item Bảng hiệu năng đã đo lại bằng \code{backend/tools/do\_hieu\_nang.py} trên
        đúng máy dùng khi bảo vệ; cấu hình máy và ngày đo khớp với báo cáo.
  \item Hệ thống chạy được trên máy khác, theo đúng các bước ở Mục A.
  \item Đã chạy thử demo trên đúng môi trường sẽ dùng khi bảo vệ.
  \item Mọi hình, bảng, giải thuật và mã nguồn đều có chú thích, có nhãn và
        được nhắc tới trong thân bài.
  \item Mọi mục tiêu ở Chương~\ref{ch:modau} đều được đối chiếu ở
        Chương~\ref{ch:ketqua}.
  \item Mọi tài liệu trong \code{refs.bib} đều được trích dẫn ít nhất một lần.
  \item Đã khai báo phần nào có dùng công cụ trí tuệ nhân tạo hỗ trợ, và giải
        thích được toàn bộ hệ thống khi bị hỏi.
\end{tukiem}

\section*{E. Nhận xét của giảng viên hướng dẫn}
\phantomsection
\addcontentsline{toc}{section}{E. Nhận xét của giảng viên hướng dẫn}

\khungnhanxet[7.5cm]

\vspace{0.8em}
\noindent\begin{minipage}[t]{0.5\textwidth}
\color{daunavy}
{\color{daumaroon}\bfseries Điểm số}

\vspace{0.7em}
\dongchamnhan{Bằng số:}

\vspace{0.9em}
\dongchamnhan{Bằng chữ:}
\end{minipage}\hfill
\okyten{\dauGVHD}{Giảng viên hướng dẫn}
